谱曲面的两处无穷远分别为 $z=\infty$ 与 $z=0$. 以下在固定 $t$ 下推导修正 Jost 解及透射系数的展开. 为使一阶系数及二阶余项均可控制, 假设
\begin{equation}\label{eq:highenergy-assumption}
q-q_0\in W^{2,1}(\mathbb R).
\end{equation}
此时 $q$ 有界, $q_x$ 连续且在两端趋于零, $q_x,q_{xx}\in L^1(\mathbb R)$.

\begin{proposition}[两处无穷远的渐近展开]\label{prop:highenergy}
在条件 \eqref{eq:highenergy-assumption} 下, 记
\begin{equation}\label{eq:asymptotic-coefficients}
B(q)=\sg Q=\begin{pmatrix}0&q\\\conj{q}&0\end{pmatrix},\qquad
J_\pm(x)=\int_{\pm\infty}^{x}(q_0^2-|q(y)|^2)\,\dd y.
\end{equation}
各 Jost 解析列在相应闭域内满足
\begin{equation}\label{eq:mu-highenergy}
\mu_\pm(x,z)=I+\dfrac{A_\pm(x)}{z}+O(z^{-2}),\qquad
A_\pm=-\ii B(q)+\ii J_\pm\sg,\qquad z\to\infty,
\end{equation}
余项对实 $x$ 一致, 矩阵式按各自解析列理解. 令
\[
J_{\rm tot}=\int_{-\infty}^{+\infty}(q_0^2-|q|^2)\,\dd x,
\]
则透射系数满足
\begin{equation}\label{eq:transmission-highenergy}
\begin{aligned}
s_{22}(z)&=1+\dfrac{\ii J_{\rm tot}}{z}+O(z^{-2}),&&z\to\infty,\quad z\in\overline{D_+},\\
s_{11}(z)&=1-\dfrac{\ii J_{\rm tot}}{z}+O(z^{-2}),&&z\to\infty,\quad z\in\overline{D_-}.
\end{aligned}
\end{equation}
在另一处无穷远,
\begin{equation}\label{eq:smallmu}
\mu_\pm(x,z)=-\dfrac{\ii q_0}{z}\sigma_1
+\dfrac{1}{q_0}\begin{pmatrix}q(x)&-J_\pm(x)\\J_\pm(x)&\conj{q(x)}\end{pmatrix}
+O(z),\qquad z\to0,
\end{equation}
以及
\begin{equation}\label{eq:transmission-smallenergy}
\begin{aligned}
s_{22}(z)&=1+\dfrac{\ii J_{\rm tot}}{q_0^2}z+O(z^2),&&z\to0,\quad z\in\overline{D_+},\\
s_{11}(z)&=1-\dfrac{\ii J_{\rm tot}}{q_0^2}z+O(z^2),&&z\to0,\quad z\in\overline{D_-}.
\end{aligned}
\end{equation}
完整散射矩阵在实轴两端满足
\begin{equation}\label{eq:scattering-two-infinities}
\begin{aligned}
S(z)&=I-\dfrac{\ii J_{\rm tot}}{z}\sg+O(z^{-2}),&&z\to\infty,\\
S(z)&=I-\dfrac{\ii J_{\rm tot}}{q_0^2}z\sg+O(z^2),&&z\to0.
\end{aligned}
\end{equation}
\end{proposition}
\begin{proof}
先确定展开的系数. 将空间方程写为
\begin{equation}\label{eq:mu-highenergy-system}
\mu_x=-\dfrac{\ii z}{2}[\sg,\mu]+Q\mu
+\dfrac{\ii q_0^2}{2z}(\sg\mu+\mu\sg).
\end{equation}
以 $\mu=I+\dfrac{A}{z}+\dfrac{A_2}{z^2}+\cdots$ 为待证展开, 比较常数项与 $z^{-1}$ 项, 分别得到
\[
-\dfrac{\ii}{2}[\sg,A]+Q=0,\qquad
A_x=-\dfrac{\ii}{2}[\sg,A_2]+QA+\ii q_0^2\sg.
\]
第一式确定 $A$ 的非对角部分为 $-\ii\sg Q=-\ii B(q)$. 对第二式取对角部分, 使用 $QB=|q|^2\sg$, 得到
\[
(A_{\rm d})_x=\ii(q_0^2-|q|^2)\sg,\qquad
A_{\rm d}(\pm\infty)=0,
\]
故 $A_\pm=-\ii B+\ii J_\pm\sg$. 以下由积分方程证明此展开及其一致余项, 无须预先假定渐近级数可逐项求导.

作局部矩阵变换
\[
 T(x,z)=I-\dfrac{\ii B(q(x))}{z},\qquad \mu_\pm= T n_\pm.
\]
由 $B\sg=-\sg B$, $B^2=|q|^2I$,
\[
\det T=1+\dfrac{|q|^2}{z^2},\qquad
 T^{-1}=\dfrac{I+\dfrac{\ii B}{z}}{1+\dfrac{|q|^2}{z^2}}.
\]
当 $|z|$ 足够大时, $ T, T^{-1}$ 对所有实 $x$ 一致有界; 在归一化端 $ T\to Y$, 故 $n_\pm\to I$. 利用 $k+\lambda=z$ 和 $\lambda-k=\dfrac{q_0^2}{z}$, 有
\begin{align*}
U T+\ii\lambda T\sg- T_x
&=\ii(\lambda-k)\sg+Q-\dfrac{k+\lambda}{z}\sg B
-\dfrac{\ii}{z}QB+\dfrac{\ii}{z}B_x\\
&=\dfrac{\ii}{z}\bigl((q_0^2-|q|^2)\sg+B_x\bigr).
\end{align*}
因此
\begin{equation}\label{eq:n-volterra-system}
(n_\pm)_x=-\ii\lambda[\sg,n_\pm]+ Cn_\pm,\qquad
 C=\dfrac{1}{z} T^{-1} C_1,\qquad
 C_1=\ii\bigl((q_0^2-|q|^2)\sg+B_x\bigr).
\end{equation}
由 \eqref{eq:highenergy-assumption}, $q_0^2-|q|^2,q_x\in L^1$, 从而
\begin{equation}\label{eq:asymptotic-L1-control}
\| C(\cdot,z)\|_{L^1}\le\dfrac{C}{|z|},\qquad
\left\| C(\cdot,z)-\dfrac{ C_1}{z}\right\|_{L^1}
\le\dfrac{C}{|z|^2}.
\end{equation}
第 $j$ 列的 Volterra 方程为
\[
n_{\pm,j}(x,z)=e_j+\int_{\pm\infty}^{x}
e^{-\ii\lambda(x-y)(\sg-s_jI)} C(y,z)n_{\pm,j}(y,z)\,\dd y.
\]
其中 $s_1=1,s_2=-1$, 传播核在相应闭解析域及积分方向上范数不超过 1. 与定理 \ref{thm:jost} 相同的有序积分估计给出第 $l$ 次迭代项的界 $\dfrac{C^l}{l!|z|^l}$, 所以
\[
\|n_\pm-I\|_{L^\infty_x}\le e^{\frac{C}{|z|}}-1=O(z^{-1}),\qquad
\sum\limits_{l=2}^{\infty}\dfrac{C^l}{l!|z|^l}=O(z^{-2}).
\]
结合 \eqref{eq:asymptotic-L1-control}, 只需计算以 $\dfrac{ C_1}{z}$ 代替 $ C$ 的第一次迭代. 以负端第一列为例, 对角分量为
\[
\dfrac{1}{z}\int_{-\infty}^{x}( C_1)_{11}(y)\,\dd y
=\dfrac{\ii J_-(x)}{z},
\]
非对角分量为
\[
\dfrac{\ii}{z}\int_{-\infty}^{x}e^{2\ii\lambda(x-y)}\conj{q_y(y)}\,\dd y.
\]
由于 $q_x$ 在负端趋于零, 分部积分得
\[
\int_{-\infty}^{x}e^{2\ii\lambda(x-y)}\conj{q_y(y)}\,\dd y
=-\dfrac{\conj{q_x(x)}}{2\ii\lambda}
+\dfrac{1}{2\ii\lambda}\int_{-\infty}^{x}e^{2\ii\lambda(x-y)}\conj{q_{yy}(y)}\,\dd y.
\]
在 $\overline{D_+}$ 中, $y\le x$ 时指数模不超过 1, 且当 $z\to\infty$ 时 $|\lambda|\ge c|z|$. 因此该非对角分量绝对值至多为
\[
\dfrac{\|q_x\|_{L^\infty}+\|q_{xx}\|_{L^1}}{2|z\lambda|}=O(z^{-2}).
\]
正端和第二列只改变积分方向或指数符号, 具有相同估计. 于是
\[
n_\pm=I+\dfrac{\ii J_\pm\sg}{z}+O(z^{-2}),
\]
与 $ T$ 相乘即证明 \eqref{eq:mu-highenergy}. 此外, 从 \eqref{eq:n-volterra-system} 及 $n_\pm$ 非对角部分为 $O(z^{-2})$, 可得 $(n_\pm)_x=O(z^{-1})$, 因而 $(\mu_\pm)_x=O(z^{-1})$, 仍按解析列理解.

将展开代入 Wronskian 表达式, 有
\[
\det(\mu_{-,1},\mu_{+,2})
=1+\dfrac{\ii(J_--J_+)}{z}+O(z^{-2}),\qquad
J_--J_+=J_{\rm tot},\qquad \gamma=1+O(z^{-2}).
\]
这给出 $s_{22}$ 的展开. 同理,
\[
\det(\mu_{+,1},\mu_{-,2})
=1-\dfrac{\ii J_{\rm tot}}{z}+O(z^{-2}),
\]
从而得到 \eqref{eq:transmission-highenergy}. 在实轴上 $\theta$ 为实数, 由散射关系
\[
\begin{aligned}
S&=e^{-\ii\theta\sg}\mu_-^{-1}\mu_+e^{\ii\theta\sg},\\
\mu_-^{-1}\mu_+&=I+\dfrac{A_+-A_-}{z}+O(z^{-2})
=I-\dfrac{\ii J_{\rm tot}}{z}\sg+O(z^{-2}),
\end{aligned}
\]
得到 \eqref{eq:scattering-two-infinities} 的第一式.

最后, 令 $w=h(z)=-\dfrac{q_0^2}{z}$. 第 \ref{ch:symmetry} 节的换叶关系给出
\[
\mu_\pm(z)=-\dfrac{\ii q_0}{z}\mu_\pm(w)\sigma_1
=-\dfrac{\ii q_0}{z}\sigma_1+\dfrac{\ii}{q_0}A_\pm\sigma_1+O(z).
\]
计算 $\ii A_\pm\sigma_1=B\sigma_1-J_\pm\sg\sigma_1$ 即得 \eqref{eq:smallmu}. 因 $h$ 交换 $D_\pm$, 透射系数满足 $s_{22}(z)=s_{11}(w)$ 和 $s_{11}(z)=s_{22}(w)$; 将 $w^{-1}=-\dfrac{z}{q_0^2}$ 代入 \eqref{eq:transmission-highenergy}, 得到 \eqref{eq:transmission-smallenergy}. 在实轴上再用 $S(z)=\sigma_1S(w)\sigma_1$, 即得 \eqref{eq:scattering-two-infinities} 的第二式.
\end{proof}

两处无穷远控制 RHP 的归一化及光场重构. 分支点 $z=\pm\ii q_0$ 的连续延拓由命题 \ref{prop:jost-branch} 给出; 轮廓交点 $z=\pm q_0$ 是不同谱弧的交汇处, 共振原像 $z_a^\pm$ 则来自传播 Lax 矩阵的有理项.


\subsubsection*{两处无穷远展开的全部矩阵元与余项计算}
\paragraph{式 \eqref{eq:highenergy-assumption} 的具体用途.}
$q-q_0\in W^{2,1}$ 蕴含 $q-q_0,q_x,q_{xx}\in L^1$.
一维绝对连续代表满足
\[
q_x(x)=-\int_x^{+\infty}q_{yy}(y)\,\dd y,
\qquad
|q_x(x)|\le\|q_{xx}\|_{L^1}.
\]
同样从负端积分也可得到有界性与两端极限为零.
$q-q_0$ 亦有界, 所以 $q$ 有界.
对原文中对角积分的被积函数, 可以精确分解为
\begin{align*}
q_0^2-|q|^2
&=-(q-q_0)\conj{q}-q_0(\conj{q}-q_0),\\
|q_0^2-|q|^2|
&\le(|q|+q_0)|q-q_0|.
\end{align*}
因此该函数属于 $L^1$,
\[
|J_\pm(x)|\le(\|q\|_{L^\infty}+q_0)\|q-q_0\|_{L^1}
\]
对所有实 $x$ 一致成立.
$J_+$ 的正端定向要特别留意:
\[
J_-(x)=\int_{-\infty}^{x}(q_0^2-|q|^2)\,\dd y,
\qquad
J_+(x)=-\int_x^{+\infty}(q_0^2-|q|^2)\,\dd y.
\]
两者的空间导数相同,
\[
J_-'(x)=J_+'(x)=q_0^2-|q(x)|^2.
\]
但两者分别在各自归一化端为零.

\paragraph{式 \eqref{eq:mu-highenergy-system} 由原空间方程逐项得到.}
直接把 $k,\lambda$ 的有理表达式代入
$\mu_x=(-\ii k\sg+Q)\mu+\ii\lambda\mu\sg$,
\begin{align*}
\mu_x
&=-\dfrac{\ii z}{2}\sg\mu
+\dfrac{\ii q_0^2}{2z}\sg\mu
+Q\mu
+\dfrac{\ii z}{2}\mu\sg
+\dfrac{\ii q_0^2}{2z}\mu\sg\\
&=-\dfrac{\ii z}{2}(\sg\mu-\mu\sg)
+Q\mu+\dfrac{\ii q_0^2}{2z}(\sg\mu+\mu\sg).
\end{align*}
对易子与反对易部分逐元为
\[
[\sg,\mu]=\begin{pmatrix}0&2\mu_{12}\\-2\mu_{21}&0\end{pmatrix},
\qquad
\sg\mu+\mu\sg=\begin{pmatrix}2\mu_{11}&0\\0&-2\mu_{22}\end{pmatrix}.
\]
而势矩阵左乘的结果是
\[
Q\mu=\begin{pmatrix}
q\mu_{21}&q\mu_{22}\\
-\conj{q}\mu_{11}&-\conj{q}\mu_{12}
\end{pmatrix}.
\]
于是完整的四个标量微分方程为
\begin{align*}
(\mu_{11})_x&=q\mu_{21}+\dfrac{\ii q_0^2}{z}\mu_{11},\\
(\mu_{12})_x&=-\ii z\mu_{12}+q\mu_{22},\\
(\mu_{21})_x&=\ii z\mu_{21}-\conj{q}\mu_{11},\\
(\mu_{22})_x&=-\conj{q}\mu_{12}-\dfrac{\ii q_0^2}{z}\mu_{22}.
\end{align*}
从这四式可以逐个核对下面每一个系数的符号.

\paragraph{式 \eqref{eq:mu-highenergy}: 常数阶首先确定两个非对角元.}
按原文的待证形式把 $\mu=I+\dfrac{A}{z}+\dfrac{A_2}{z^2}+\cdots$ 代入.
左侧没有常数项, 右侧常数阶为
\[
-\dfrac{\ii}{2}[\sg,A]+Q
=\begin{pmatrix}
0&-\ii A_{12}+q\\
\ii A_{21}-\conj{q}&0
\end{pmatrix}=0.
\]
两个非平凡的标量等式解为
\[
-\ii A_{12}+q=0\ \Longrightarrow\ A_{12}=-\ii q,
\qquad
\ii A_{21}-\conj{q}=0\ \Longrightarrow\ A_{21}=-\ii\conj{q}.
\]
这个常数阶方程没有确定 $A_{11},A_{22}$,
因为与 $\sg$ 的对易子对角元恒为零.
不能把这两个尚未确定的对角系数直接设为零.

$z^{-1}$ 阶的矩阵方程中,
\[
-\dfrac{\ii}{2}[\sg,A_2]
=\begin{pmatrix}0&-\ii(A_2)_{12}\\
\ii(A_2)_{21}&0\end{pmatrix},
\qquad
QA=\begin{pmatrix}
qA_{21}&qA_{22}\\-\conj{q}A_{11}&-\conj{q}A_{12}
\end{pmatrix}.
\]
逐个比较四个矩阵元,
\begin{align*}
(A_{11})_x&=qA_{21}+\ii q_0^2,\\
(A_{12})_x&=-\ii(A_2)_{12}+qA_{22},\\
(A_{21})_x&=\ii(A_2)_{21}-\conj{q}A_{11},\\
(A_{22})_x&=-\conj{q}A_{12}-\ii q_0^2.
\end{align*}
将已经算出的 $A_{12},A_{21}$ 代入两条对角方程,
\begin{align*}
(A_{11})_x
&=q(-\ii\conj{q})+\ii q_0^2
=\ii(q_0^2-|q|^2),\\
(A_{22})_x
&=-\conj{q}(-\ii q)-\ii q_0^2
=-\ii(q_0^2-|q|^2).
\end{align*}
在归一化端 $\mu_\pm\to Y=I-\dfrac{\ii q_0}{z}\sigma_1$,
所以对角一阶系数的端态为零.
从相应端积分,
\begin{align*}
(A_\pm)_{11}(x)
&=\ii\int_{\pm\infty}^{x}(q_0^2-|q(y)|^2)\,\dd y
=\ii J_\pm(x),\\
(A_\pm)_{22}(x)
&=-\ii\int_{\pm\infty}^{x}(q_0^2-|q(y)|^2)\,\dd y
=-\ii J_\pm(x).
\end{align*}
因此四个系数完整合起来为
\[
A_\pm=\begin{pmatrix}
\ii J_\pm&-\ii q\\-\ii\conj{q}&-\ii J_\pm
\end{pmatrix}
=-\ii B(q)+\ii J_\pm\sg.
\]
这是系数的形式计算; 原文随后用积分方程证明余项,
所以不需要先假定整个渐近级数可以无限次逐项求导.

\paragraph{继续核对下一阶的两个非对角系数.}
虽然原结论只需一阶系数, 上面四个矩阵元还显示形式上的
$(A_2)_{12},(A_2)_{21}$ 应满足什么关系.
用 $A_{12}=-\ii q,A_{22}=-\ii J_\pm$,
\[
-\ii q_x=-\ii(A_2)_{12}-\ii qJ_\pm,
\qquad (A_2)_{12}=q_x-qJ_\pm.
\]
用 $A_{21}=-\ii\conj{q},A_{11}=\ii J_\pm$,
\[
-\ii\conj{q_x}=\ii(A_2)_{21}-\ii\conj{q}J_\pm,
\qquad (A_2)_{21}=-\conj{q_x}+\conj{q}J_\pm.
\]
这只是对待证级数的代数核对, 并不把原文的
$O(z^{-2})$ 余项擅自改成已证明的完整二阶展开.
原论证不需要对 $A_2$ 的对角部分再作积分.

\paragraph{原文局部矩阵 $ T$ 的逆逐元验证.}
这里仅使用原文已经引入的 $ T,n_\pm, C, C_1$,
不新增任何为简算而设的变换或辅助未知量.
\[
B=\begin{pmatrix}0&q\\\conj{q}&0\end{pmatrix},
\qquad
B^2=\begin{pmatrix}|q|^2&0\\0&|q|^2\end{pmatrix}=|q|^2I.
\]
于是
\[
 T=\begin{pmatrix}
1&-\dfrac{\ii q}{z}\\-\dfrac{\ii\conj{q}}{z}&1
\end{pmatrix},
\qquad
\det T
=1-\left(-\dfrac{\ii q}{z}\right)
\left(-\dfrac{\ii\conj{q}}{z}\right)
=1+\dfrac{|q|^2}{z^2}.
\]
将所写的逆矩阵与原矩阵相乘,
\begin{align*}
\left(I-\dfrac{\ii B}{z}\right)
\left(I+\dfrac{\ii B}{z}\right)
&=I+\dfrac{B^2}{z^2}
=\left(1+\dfrac{|q|^2}{z^2}\right)I.
\end{align*}
这说明逆矩阵公式中分母的符号为正.
只要 $|z|^2\ge2\|q\|_{L^\infty}^2$, 就有
\[
\left|1+\dfrac{|q(x)|^2}{z^2}\right|
\ge1-\dfrac{|q(x)|^2}{|z|^2}\ge\dfrac{1}{2}
\]
对所有 $x$ 一致成立, 因而逆矩阵一致有界.
在归一化端 $q\to q_0$, $ T\to Y$,
所以 $n_\pm= T^{-1}\mu_\pm\to I$.

\paragraph{式 \eqref{eq:n-volterra-system} 的剩余矩阵逐个矩阵元相消.}
先把 $U T$ 直接乘出来:
\[
U T=\begin{pmatrix}
-\ii k-\dfrac{\ii|q|^2}{z}&q-\dfrac{kq}{z}\\
-\conj{q}+\dfrac{k\conj{q}}{z}&\ii k+\dfrac{\ii|q|^2}{z}
\end{pmatrix}.
\]
另外两个矩阵分别为
\[
\ii\lambda T\sg
=\begin{pmatrix}
\ii\lambda&-\dfrac{\lambda q}{z}\\
\dfrac{\lambda\conj{q}}{z}&-\ii\lambda
\end{pmatrix},
\qquad
- T_x
=\begin{pmatrix}0&\dfrac{\ii q_x}{z}\\
\dfrac{\ii\conj{q_x}}{z}&0\end{pmatrix}.
\]
逐元相加, 左上元是
\[
-\ii k-\dfrac{\ii|q|^2}{z}+\ii\lambda
=\ii(\lambda-k)-\dfrac{\ii|q|^2}{z}
=\dfrac{\ii(q_0^2-|q|^2)}{z}.
\]
右下元是它的相反数.
右上元为
\[
q-\dfrac{(k+\lambda)q}{z}+\dfrac{\ii q_x}{z}
=q-q+\dfrac{\ii q_x}{z}
=\dfrac{\ii q_x}{z}.
\]
左下元为
\[
-\conj{q}+\dfrac{(k+\lambda)\conj{q}}{z}
+\dfrac{\ii\conj{q_x}}{z}
=-\conj{q}+\conj{q}+\dfrac{\ii\conj{q_x}}{z}
=\dfrac{\ii\conj{q_x}}{z}.
\]
因此剩余矩阵恰好为
\[
U T+\ii\lambda T\sg- T_x
=\dfrac{ C_1}{z},
\qquad
 C_1=\begin{pmatrix}
\ii(q_0^2-|q|^2)&\ii q_x\\
\ii\conj{q_x}&-\ii(q_0^2-|q|^2)
\end{pmatrix}.
\]

\paragraph{为什么 $n$ 方程中还有一个对易子.}
将 $\mu= Tn$ 代入原空间方程,
\[
 T_xn+ Tn_x
=U Tn+\ii\lambda Tn\sg.
\]
不能把最后一项直接改成 $\ii\lambda T\sg n$.
应加减同一项, 按原顺序写成
\begin{align*}
 Tn_x
&=(U T+\ii\lambda T\sg- T_x)n
+\ii\lambda T(n\sg-\sg n)\\
&=\dfrac{ C_1}{z}n
-\ii\lambda T[\sg,n].
\end{align*}
左乘 $ T^{-1}$ 就得到
$n_x=-\ii\lambda[\sg,n]+ Cn$.
整个推导中没有把非交换的矩阵乘法次序互换.

\paragraph{式 \eqref{eq:asymptotic-L1-control} 的两条估计.}
其中平方模差的可积性可逐项估计为
\begin{align*}
\left|q_0^2-|q|^2\right|
&=(q_0+|q|)\left|q_0-|q|\right|\\
&\le(q_0+\|q\|_{L^\infty})|q-q_0|,\\
\int_{-\infty}^{+\infty}\left|q_0^2-|q(x)|^2\right|\,\dd x
&\le(q_0+\|q\|_{L^\infty})\|q-q_0\|_{L^1}<\infty.
\end{align*}
第一行是平方差因式分解, 第二行使用反三角不等式.
由被积函数的可积性,
\[
\| C_1\|_{L^1}
\le C\left(\|q_0^2-|q|^2\|_{L^1}+\|q_x\|_{L^1}\right)<\infty.
\]
这里用于不等式上界的 $C>0$ 是与大谱参数无关的标量估计常数. 它与左侧的矩阵函数 $C(x,z)$ 及其系数 $C_1(x)$ 由自变量和维数区分. 第7部分的 $T(x,z)$ 同样是二阶矩阵, 与第8部分的标量指数 $T=e^{2s}$ 分别使用各自的定义.
原文的逆矩阵可改写成精确差值
\[
 T^{-1}-I
=\dfrac{\dfrac{\ii B}{z}-\dfrac{|q|^2}{z^2}I}
{1+\dfrac{|q|^2}{z^2}}.
\]
所以在前面的大 $|z|$ 范围,
\[
\| T^{-1}-I\|_{L^\infty_x}\le\dfrac{C}{|z|}.
\]
再按矩阵乘法范数估计,
\begin{align*}
\| C\|_{L^1}
&=\left\|\dfrac{1}{z} T^{-1} C_1\right\|_{L^1}
\le\dfrac{\| T^{-1}\|_{L^\infty}\| C_1\|_{L^1}}{|z|}
\le\dfrac{C}{|z|},\\
\left\| C-\dfrac{ C_1}{z}\right\|_{L^1}
&=\left\|\dfrac{1}{z}( T^{-1}-I) C_1\right\|_{L^1}\\
&\le\dfrac{\| T^{-1}-I\|_{L^\infty}\| C_1\|_{L^1}}{|z|}
\le\dfrac{C}{|z|^2}.
\end{align*}
因此不是把 $ T^{-1}$ 形式上替换为 $I$,
而是有一个可积的 $O(z^{-2})$ 误差.

\paragraph{四个 $n$ 列的传播核及第一次迭代.}
第一列对应
\[
e^{-\ii\lambda(x-y)(\sg-I)}
=\begin{pmatrix}1&0\\0&e^{2\ii\lambda(x-y)}\end{pmatrix},
\]
第二列对应
\[
e^{-\ii\lambda(x-y)(\sg+I)}
=\begin{pmatrix}e^{-2\ii\lambda(x-y)}&0\\0&1\end{pmatrix}.
\]
四种列与允许方向与第3部分相同:
负端第一列和正端第二列用 $\operatorname{Im}\lambda\ge0$,
正端第一列和负端第二列用 $\operatorname{Im}\lambda\le0$.
故每个传播对角元的模不超过 1.
把 $\dfrac{ C_1}{z}$ 作用在自由项 $e_1,e_2$ 上,
第一次迭代分别为
\begin{align*}
n_{\pm,1}(x,z)
&=e_1+\dfrac{\ii}{z}
\begin{pmatrix}
\displaystyle\int_{\pm\infty}^{x}(q_0^2-|q|^2)\,\dd y\\[4pt]
\displaystyle\int_{\pm\infty}^{x}
 e^{2\ii\lambda(x-y)}\conj{q_y(y)}\,\dd y
\end{pmatrix}+O(z^{-2}),\\
n_{\pm,2}(x,z)
&=e_2+\dfrac{\ii}{z}
\begin{pmatrix}
\displaystyle\int_{\pm\infty}^{x}
 e^{-2\ii\lambda(x-y)}q_y(y)\,\dd y\\[4pt]
-\displaystyle\int_{\pm\infty}^{x}(q_0^2-|q|^2)\,\dd y
\end{pmatrix}+O(z^{-2}).
\end{align*}
这里的 $O(z^{-2})$ 包含系数替换误差与第二次以后的迭代.
具体地, 第 $l$ 次有序积分的上界为
$\dfrac{C^l}{l!|z|^l}$.
当 $|z|\ge1$ 时,
\begin{align*}
\sum\limits_{l=2}^{\infty}\dfrac{C^l}{l!|z|^l}
&\le\dfrac{1}{|z|^2}\sum\limits_{l=2}^{\infty}\dfrac{C^l}{l!}
=\dfrac{e^C-1-C}{|z|^2}.
\end{align*}
因此从第一次迭代开始分离余项是有定量依据的.

\paragraph{分部积分给出的额外一个大谱参数因子.}
对第一列的非对角积分,
\[
\partial_y e^{2\ii\lambda(x-y)}
=-2\ii\lambda e^{2\ii\lambda(x-y)}.
\]
把指数写成它的 $y$ 导数乘系数, 再分部积分,
\begin{align*}
&\int_{\pm\infty}^{x}e^{2\ii\lambda(x-y)}\conj{q_y(y)}\,\dd y\\
&=-\dfrac{1}{2\ii\lambda}
\int_{\pm\infty}^{x}\partial_y\bigl(e^{2\ii\lambda(x-y)}\bigr)
\conj{q_y(y)}\,\dd y\\
&=-\dfrac{1}{2\ii\lambda}
\left[e^{2\ii\lambda(x-y)}\conj{q_y(y)}\right]_{\pm\infty}^{x}
+\dfrac{1}{2\ii\lambda}
\int_{\pm\infty}^{x}e^{2\ii\lambda(x-y)}\conj{q_{yy}(y)}\,\dd y\\
&=-\dfrac{\conj{q_x(x)}}{2\ii\lambda}
+\dfrac{1}{2\ii\lambda}
\int_{\pm\infty}^{x}e^{2\ii\lambda(x-y)}\conj{q_{yy}(y)}\,\dd y.
\end{align*}
无穷远的边界项为零, 因为指数在当前允许方向有界且 $q_x\to0$.
第二列的指数导数符号相反,
\[
\partial_y e^{-2\ii\lambda(x-y)}
=2\ii\lambda e^{-2\ii\lambda(x-y)},
\]
所以逐项计算为
\begin{align*}
&\int_{\pm\infty}^{x}e^{-2\ii\lambda(x-y)}q_y(y)\,\dd y\\
&=\dfrac{1}{2\ii\lambda}
\left[e^{-2\ii\lambda(x-y)}q_y(y)\right]_{\pm\infty}^{x}
-\dfrac{1}{2\ii\lambda}
\int_{\pm\infty}^{x}e^{-2\ii\lambda(x-y)}q_{yy}(y)\,\dd y\\
&=\dfrac{q_x(x)}{2\ii\lambda}
-\dfrac{1}{2\ii\lambda}
\int_{\pm\infty}^{x}e^{-2\ii\lambda(x-y)}q_{yy}(y)\,\dd y.
\end{align*}
正端积分的方向仍为从 $+\infty$ 到 $x$;
定向积分的分部积分公式本身不需要另加符号.
两类积分的绝对值都不超过
\[
\dfrac{\|q_x\|_{L^\infty}+\|q_{xx}\|_{L^1}}{2|\lambda|}.
\]
对大 $|z|$,
\[
|\lambda(z)|
=\dfrac{1}{2}\left|z+\dfrac{q_0^2}{z}\right|
\ge\dfrac{1}{2}\left(|z|-\dfrac{q_0^2}{|z|}\right)
\ge\dfrac{|z|}{4}
\]
只需 $|z|^2\ge2q_0^2$.
所以第一次迭代前面的 $\dfrac{1}{z}$ 再乘这里的 $\dfrac{1}{\lambda}$,
确实给出 $O(z^{-2})$.

\paragraph{乘回 $ T$ 时四个矩阵元分别是什么.}
由对角积分与刚才的非对角估计,
\[
n_\pm=\begin{pmatrix}
1+\dfrac{\ii J_\pm}{z}+O(z^{-2})&O(z^{-2})\\
O(z^{-2})&1-\dfrac{\ii J_\pm}{z}+O(z^{-2})
\end{pmatrix}.
\]
直接把 $\mu_\pm= Tn_\pm$ 逐元相乘,
\begin{align*}
(\mu_\pm)_{11}
&=(n_\pm)_{11}-\dfrac{\ii q}{z}(n_\pm)_{21}
=1+\dfrac{\ii J_\pm}{z}+O(z^{-2}),\\
(\mu_\pm)_{12}
&=(n_\pm)_{12}-\dfrac{\ii q}{z}(n_\pm)_{22}
=-\dfrac{\ii q}{z}+O(z^{-2}),\\
(\mu_\pm)_{21}
&=-\dfrac{\ii\conj{q}}{z}(n_\pm)_{11}+(n_\pm)_{21}
=-\dfrac{\ii\conj{q}}{z}+O(z^{-2}),\\
(\mu_\pm)_{22}
&=-\dfrac{\ii\conj{q}}{z}(n_\pm)_{12}+(n_\pm)_{22}
=1-\dfrac{\ii J_\pm}{z}+O(z^{-2}).
\end{align*}
这里 $q,J_\pm$ 有界, 所以余项乘上这些系数仍有同一个阶.
四个余项在各自解析列允许的闭域与所有实 $x$ 上一致.

空间导数的估计也可以从 $n$ 方程逐元核对:
\[
[\sg,n_\pm]=\begin{pmatrix}0&O(z^{-2})\\O(z^{-2})&0\end{pmatrix}.
\]
因此 $\lambda[\sg,n_\pm]=O(z^{-1})$,
$ Cn_\pm=O(z^{-1})$, 故 $(n_\pm)_x=O(z^{-1})$.
同时 $ T_x=O(z^{-1})$, 所以
\[
(\mu_\pm)_x= T_xn_\pm+ T(n_\pm)_x=O(z^{-1}).
\]
这一步不是直接对没有先验证的 $O(z^{-2})$ 余项求导.

\paragraph{式 \eqref{eq:transmission-highenergy}: 两个 Wronskian 都展开.}
由四个矩阵元,
\begin{align*}
\det(\mu_{-,1},\mu_{+,2})
&=(\mu_-)_{11}(\mu_+)_{22}-(\mu_+)_{12}(\mu_-)_{21}\\
&=\left(1+\dfrac{\ii J_-}{z}+O(z^{-2})\right)
\left(1-\dfrac{\ii J_+}{z}+O(z^{-2})\right)\\
&\quad-\left(-\dfrac{\ii q}{z}+O(z^{-2})\right)
\left(-\dfrac{\ii\conj{q}}{z}+O(z^{-2})\right)\\
&=1+\dfrac{\ii(J_--J_+)}{z}+O(z^{-2}).
\end{align*}
非对角乘积至少是二阶, 所以不会改变一阶系数.
两个定向积分的差为
\begin{align*}
J_--J_+
&=\int_{-\infty}^{x}(q_0^2-|q|^2)\,\dd y
+\int_x^{+\infty}(q_0^2-|q|^2)\,\dd y\\
&=\int_{-\infty}^{+\infty}(q_0^2-|q|^2)\,\dd y=J_{\rm tot}.
\end{align*}
而
\[
\dfrac{1}{\gamma(z)}
=\dfrac{1}{1+\dfrac{q_0^2}{z^2}}
=1-\dfrac{q_0^2}{z^2}+O(z^{-4})
\]
没有一阶项, 所以 $s_{22}=1+\dfrac{\ii J_{\rm tot}}{z}+O(z^{-2})$.
另一个 Wronskian逐项为
\begin{align*}
\det(\mu_{+,1},\mu_{-,2})
&=(\mu_+)_{11}(\mu_-)_{22}-(\mu_-)_{12}(\mu_+)_{21}\\
&=\left(1+\dfrac{\ii J_+}{z}+O(z^{-2})\right)
\left(1-\dfrac{\ii J_-}{z}+O(z^{-2})\right)\\
&\quad-\left(-\dfrac{\ii q}{z}+O(z^{-2})\right)
\left(-\dfrac{\ii\conj{q}}{z}+O(z^{-2})\right)\\
&=1+\dfrac{\ii(J_+-J_-)}{z}+O(z^{-2})\\
&=1-\dfrac{\ii J_{\rm tot}}{z}+O(z^{-2}).
\end{align*}
除以同一个 $\gamma$ 就得到 $s_{11}$ 的展开.

\paragraph{两个非对角散射系数为什么没有一阶项.}
在实轴上, 两个 Wronskian 的一阶项直接抵消:
\begin{align*}
\det(\mu_{-,1},\mu_{+,1})
&=(\mu_-)_{11}(\mu_+)_{21}-(\mu_+)_{11}(\mu_-)_{21}\\
&=\left(1+\dfrac{\ii J_-}{z}\right)
\left(-\dfrac{\ii\conj{q}}{z}\right)
-\left(1+\dfrac{\ii J_+}{z}\right)
\left(-\dfrac{\ii\conj{q}}{z}\right)+O(z^{-2})\\
&=-\dfrac{\ii\conj{q}}{z}+\dfrac{\ii\conj{q}}{z}+O(z^{-2})
=O(z^{-2}),\\
\det(\mu_{+,2},\mu_{-,2})
&=(\mu_+)_{12}(\mu_-)_{22}-(\mu_-)_{12}(\mu_+)_{22}\\
&=\left(-\dfrac{\ii q}{z}\right)
\left(1-\dfrac{\ii J_-}{z}\right)
-\left(-\dfrac{\ii q}{z}\right)
\left(1-\dfrac{\ii J_+}{z}\right)+O(z^{-2})\\
&=-\dfrac{\ii q}{z}+\dfrac{\ii q}{z}+O(z^{-2})=O(z^{-2}).
\end{align*}
乘回 $e^{\pm2\ii\theta}$ 后, 在实轴上指数模为 1,
所以 $s_{12},s_{21}=O(z^{-2})$.
这说明完整散射矩阵的一阶系数是对角矩阵.
在非实谱参数上, 不能不加条件地把指数因子的模也当作 1;
这正是原文把完整矩阵的展开限制在实轴上的原因.

\paragraph{式 \eqref{eq:scattering-two-infinities} 的矩阵计算交叉核对.}
由 $\det\mu_-=\gamma$, 二阶逆矩阵公式给出
\[
\mu_-^{-1}
=\dfrac{1}{\gamma}
\begin{pmatrix}(\mu_-)_{22}&-(\mu_-)_{12}\\
-(\mu_-)_{21}&(\mu_-)_{11}\end{pmatrix}
=I-\dfrac{A_-}{z}+O(z^{-2}).
\]
因而
\begin{align*}
\mu_-^{-1}\mu_+
&=\left(I-\dfrac{A_-}{z}+O(z^{-2})\right)
\left(I+\dfrac{A_+}{z}+O(z^{-2})\right)\\
&=I+\dfrac{A_+-A_-}{z}+O(z^{-2}),\\
A_+-A_-
&=-\ii B+\ii J_+\sg+\ii B-\ii J_-\sg\\
&=\ii(J_+-J_-)\sg=-\ii J_{\rm tot}\sg.
\end{align*}
对角系数与 $e^{\ii\theta\sg}$ 交换,
所以左右相位共轭不改变它.
实轴上相位与逆相位有界, 因而也保留余项阶数,
最终得到原文 $S$ 在大 $z$ 的矩阵展开.

\paragraph{式 \eqref{eq:smallmu}: 用原换叶关系逐项乘出常数矩阵.}
这里的 $w=h(z)$ 是原文已经使用的谱点记号.
由于 $h(h(z))=z$, 在 $w$ 处应用换叶关系可得
\[
\mu_\pm(z)=\mu_\pm(w)C(w),
\qquad
C(w)=\dfrac{\ii w}{q_0}\sigma_1
=-\dfrac{\ii q_0}{z}\sigma_1.
\]
因此
\begin{align*}
\mu_\pm(z)
&=-\dfrac{\ii q_0}{z}
\left(I+\dfrac{A_\pm}{w}+O(w^{-2})\right)\sigma_1\\
&=-\dfrac{\ii q_0}{z}\sigma_1
+\dfrac{\ii}{q_0}A_\pm\sigma_1+O(z).
\end{align*}
这里使用 $w^{-1}=-\dfrac{z}{q_0^2}$;
大 $w$ 的 $O(w^{-2})$ 乘上 $\dfrac{1}{z}$ 后是 $O(z)$.
接着不省略常数矩阵乘法:
\[
B\sigma_1=\begin{pmatrix}q&0\\0&\conj{q}\end{pmatrix},
\qquad
\sg\sigma_1=\begin{pmatrix}0&1\\-1&0\end{pmatrix},
\]
\begin{align*}
\ii A_\pm\sigma_1
&=\ii(-\ii B+\ii J_\pm\sg)\sigma_1\\
&=B\sigma_1-J_\pm\sg\sigma_1\\
&=\begin{pmatrix}q&-J_\pm\\J_\pm&\conj{q}\end{pmatrix}.
\end{align*}
所以小 $z$ 的四个矩阵元分别为
\begin{align*}
(\mu_\pm)_{11}&=\dfrac{q}{q_0}+O(z),&
(\mu_\pm)_{12}&=-\dfrac{\ii q_0}{z}-\dfrac{J_\pm}{q_0}+O(z),\\
(\mu_\pm)_{21}&=-\dfrac{\ii q_0}{z}+\dfrac{J_\pm}{q_0}+O(z),&
(\mu_\pm)_{22}&=\dfrac{\conj{q}}{q_0}+O(z).
\end{align*}
非对角常数项的符号相反, 正是上面 $\sg\sigma_1$ 的两项符号不同.
在背景 $q=q_0,J_\pm=0$ 上, 这个式子恢复
$Y=I-\dfrac{\ii q_0}{z}\sigma_1$, 可直接复核常数矩阵.

\paragraph{式 \eqref{eq:transmission-smallenergy} 的每个符号.}
$h$ 交换两个解析域, 因而在 $z\in D_+$ 时 $w\in D_-$.
使用相应域内正确的透射展开,
\begin{align*}
s_{22}(z)
&=s_{11}(w)
=1-\dfrac{\ii J_{\rm tot}}{w}+O(w^{-2})\\
&=1-\ii J_{\rm tot}\left(-\dfrac{z}{q_0^2}\right)+O(z^2)
=1+\dfrac{\ii J_{\rm tot}}{q_0^2}z+O(z^2).
\end{align*}
在 $z\in D_-$ 时 $w\in D_+$, 同样有
\begin{align*}
s_{11}(z)
&=s_{22}(w)
=1+\dfrac{\ii J_{\rm tot}}{w}+O(w^{-2})\\
&=1+\ii J_{\rm tot}\left(-\dfrac{z}{q_0^2}\right)+O(z^2)
=1-\dfrac{\ii J_{\rm tot}}{q_0^2}z+O(z^2).
\end{align*}
小 $z$ 展开的正负号不仅由 $w^{-1}$ 的负号决定,
还使用了换叶后必须交换 $s_{11}$ 与 $s_{22}$.

完整散射矩阵在实轴上满足
\begin{align*}
S(z)
&=\sigma_1S(w)\sigma_1\\
&=\sigma_1\left(I-\dfrac{\ii J_{\rm tot}}{w}\sg
+O(w^{-2})\right)\sigma_1\\
&=I+\dfrac{\ii J_{\rm tot}}{w}\sg+O(w^{-2})\\
&=I-\dfrac{\ii J_{\rm tot}}{q_0^2}z\sg+O(z^2).
\end{align*}
第三行使用 $\sigma_1\sg\sigma_1=-\sg$,
第四行再使用 $w^{-1}=-\dfrac{z}{q_0^2}$.
两个非对角元也由换叶关系保持为 $O(z^2)$.
所有小 $z$ 公式都是从已验证的大 $z$ 展开与原约化关系推出,
没有在 $z=0$ 处对含奇系数的微分方程作不受控的形式除法.
